% Includable section. Requires amsmath/amssymb/amsthm and theorem,
% lemma, proposition, and remark environments in the main document.
% The conjunction ladder is existing work, not a new result.
\section{An exponential lower bound for exact circuit minimization}
\label{sec:exact-occam-lower}

The conjunction ladder below is the construction from Appendix~K of
the Vaintrob--H\"anni slides \cite{slides}, proved for unrestricted
binary-gate circuits by H\"anni and Claude \cite{soi}.
We give an elementary support-based proof and make the prediction
protocol explicit. Its lower bound is exponential in the target's
\emph{minimum gate count}, and every hypothesis used by the learner
really is a globally smallest consistent circuit.

Throughout this section a circuit is a finite, acyclic, single-output
circuit over either all sixteen two-input Boolean gates or the standard
basis consisting of binary AND and OR and unary NOT. Constants are
available as source wires, and arbitrary fan-out is allowed. The output
must be a gate, so a circuit has at least one gate. Size counts every
gate, including NOT, and $C(f)$ denotes the minimum size of a circuit
computing $f$ on the full Boolean cube.

A \emph{history-only minimum-circuit learner} chooses a circuit $h_t$
of minimum gate count among all circuits consistent with the first
$t-1$ labeled examples, and predicts $h_t(x_t)$. Its choice of $h_t$
can depend on the entire past, but not on the current query $x_t$.
For deterministic learners this includes every fixed ordering of
equal-size circuit descriptions or equal-complexity functions.

\begin{lemma}[Support forces exact minimality]
\label{lem:lower-support}
A circuit of $s$ gates depends on at most $s+1$ input variables.
Suppose $h(1^d)=1$ and $h(x)=0$ whenever $|x|<k$.
Then $h$ depends on at least $k$ variables. If it depends on exactly
$k$ variables, it is the conjunction of those variables.
\end{lemma}
\begin{proof}
Prune gates having no path to the output, leaving $t\le s$ gates.
Of their at most $2t$ input slots, at least $t-1$ reference other
remaining gates. Thus at most $t+1\le s+1$ slots can hold distinct
input variables. This argument permits unary gates and shared wires.

Let $S$ be the essential-variable set of $h$, and let $1_S$ denote
its indicator in $\{0,1\}^d$. If $|S|<k$, then
$h(1_S)=h(1^d)=1$, contradicting the lower-weight negative labels.
If $|S|=k$, extend any assignment on $S$ other than $1_S$ by zeros
outside $S$. The resulting input has weight below $k$, so its label
is zero. The all-ones assignment on $S$ has label one. Hence
$h(x)=\bigwedge_{j\in S}x_j$.
\end{proof}

\begin{theorem}[Conjunction ladder for exact minimum circuits]
\label{thm:exact-occam-ladder}
Let $d\ge3$ and $f(x)=\bigwedge_{j=1}^d x_j$.
Then $C(f)=d-1$. For every deterministic history-only
minimum-circuit learner, there is a permutation of $\{0,1\}^d$
on which its total number of mistakes is at least
\[
  \sum_{k=2}^{d-1}\binom dk
  =2^d-d-2
  =2^{C(f)+1}-C(f)-3.
\]
The permutation may depend on the learner's tie rule.
\end{theorem}
\begin{proof}
First reveal $1^d$ with label one, then $0^d$ and all singleton
inputs with label zero. For $k=2,\ldots,d-1$, process the inputs of
weight $k$, choosing the order within the level as follows.

At the beginning of this phase, every input of weight less than $k$
has been revealed with label zero. At an intermediate step, let
$\mathcal R$ be the collection of $k$-sets whose indicator inputs
have not yet been queried. If $\mathcal R\ne\varnothing$, each
\[
  h_S(x)=\bigwedge_{j\in S}x_j,\qquad S\in\mathcal R,
\]
is consistent and has a circuit of $k-1$ gates. Every consistent
circuit has at least $k-1$ gates by Lemma~\ref{lem:lower-support}.
Consequently the \emph{global} minimum over all consistent circuits
is exactly $k-1$.

Every minimizing circuit depends on at most $k$ variables and at
least $k$ variables. The same lemma therefore says that it computes
$h_S$ for a $k$-set $S$. Consistency forces $S\in\mathcal R$.
Query $1_S$. The learner predicts one and the target label is zero.
Moreover, for another $k$-set $T$,
\[
  h_T(1_S)=1\quad\Longleftrightarrow\quad T\subseteq S
  \quad\Longleftrightarrow\quad T=S.
\]
Thus this query eliminates precisely one of the surviving
$k$-conjunctions. Repeat until the level is exhausted. Every input
in every level $k=2,\ldots,d-1$ is a mistake.

The query can be determined before prediction because the learner's
current hypothesis is a deterministic function of its past.
Equivalently, simulate the construction in advance to obtain a
fixed permutation for this learner. Finally, $f$ depends on all
$d$ variables, so Lemma~\ref{lem:lower-support} gives
$C(f)\ge d-1$; the usual AND tree attains equality.
\end{proof}

\paragraph{A fixed tie rule and a fixed input order.}
There is also a direct nonadaptive example. At each size $s\ge1$,
rank canonical left-associated conjunction circuits on $(s+1)$-sets
first, ordered lexicographically by their supports, and rank all
other circuits afterward in any fixed order. Choose the first
consistent circuit of minimum size. Present $1^d$, $0^d$, all
singletons, and then the levels of weight $2,\ldots,d-1$ in that
same lexicographic order. At every step after the singleton level,
the current hypothesis is the conjunction indexed by the next
input. The number of mistakes is exactly $2^d-d-2$. Both orderings
are fixed independently of the run. The theorem is stronger in a
different direction: every deterministic history-only tie rule
has some bad ordering, without asserting one ordering bad for
all tie rules.

\paragraph{Why the singleton level is omitted.}
Under the mandatory-output-gate convention, a one-gate
$x_i\wedge x_j$ is consistent with the initial positive example,
the zero vector, and every singleton negative. Therefore the
first level need not incur any mistakes. If instead an input or
constant wire may itself be a zero-gate output, the proof starts
at $k=1$ and yields $2^d-2$ mistakes after the two initial examples,
with the same target complexity $d-1$. This stronger count uses
a different, explicitly stated circuit convention.

\paragraph{The complexity measure matters.}
The theorem is literally exponential in the minimum number of
gates. For the explicit gate-list language that encodes size in
unary and then gate types and wire addresses, the target has
description length $b=\Theta(d\log d)$. In that measure the
conclusion is $2^{\Omega(b/\log b)}$, not a proved
$2^{\Omega(b)}$ bound. A constant-factor compilation into a
NAND-only basis does not automatically transfer the theorem:
compilation need not preserve which consistent circuits are
smallest. The next section proves the NAND-only statement directly.

\paragraph{Scope of the tie-breaking quantifier.}
The deterministic theorem does not permit an adversary to inspect
a hidden fresh random circuit before selecting its query. Nor does
it cover a deterministic learner that uses the current query to
choose between minimum circuits. For example, after receiving $x_t$,
a learner may take the majority label among all current minimum
circuits and then select a minimum circuit giving that label.
This still selects a smallest consistent circuit at every round.
Each mistake halves the candidates at the current size, or exhausts
that size. If $N_s$ bounds the number of size-$s$ descriptions, its
mistakes are at most
\[
  \sum_{s=1}^{C(f)}\bigl(1+\log_2 N_s\bigr)
  =O\!\left(C(f)^2\log(d+C(f))\right).
\]
Thus specifying only ``a smallest circuit each round'' leaves a
material ambiguity: the exponential lower bound concerns a single
deterministic hypothesis chosen from the past, before using the
current query.
